\documentclass[12pt,a4paper]{article}

\usepackage[a4paper,margin=2.0cm]{geometry}
\usepackage{fontspec}
\usepackage{polyglossia}
\usepackage{enumitem}
\usepackage{hyperref}
\usepackage{longtable}
\usepackage{array}

\setmainlanguage{russian}
\setotherlanguage{english}

\setmainfont[
  Path=../fonts/,
  UprightFont=pt-astra-serif_regular.ttf,
  BoldFont=pt-astra-serif_bold.ttf,
  ItalicFont=pt-astra-serif_italic.ttf,
  BoldItalicFont=pt-astra-serif_bold-italic.ttf
]{PT Astra Serif}

\setmonofont[
  Path=../fonts/fira-code/,
  UprightFont=FiraCode-Regular.ttf,
  BoldFont=FiraCode-Bold.ttf
]{Fira Code}

\hypersetup{
  hidelinks,
  pdfauthor={Нахатакян Артур Романович},
  pdftitle={Раздаточный материал к ВКР}
}

\setlist[itemize]{leftmargin=1.4cm}
\setlist[enumerate]{leftmargin=1.4cm}
\setlength{\parindent}{0pt}
\setlength{\parskip}{0.45em}
\setlength{\tabcolsep}{6pt}
\renewcommand{\arraystretch}{1.2}
\emergencystretch=2em

\begin{document}

\begin{center}
{\Large \textbf{Раздаточный материал к выпускной квалификационной работе}}\\[0.8em]
{\normalsize \textbf{Тема:} «Разработка системы управления контентом в сфере маркетинга на базе Strapi и Astro»}\\[0.4em]
{\normalsize \textbf{Обучающийся:} Нахатакян Артур Романович}\\
{\normalsize \textbf{Научный руководитель:} Государев Илья Борисович, к.п.н., доцент}\\
{\normalsize Санкт-Петербург, 2026}
\end{center}

\section*{1. Актуальность, объект и предмет исследования}

Актуальность темы обусловлена тем, что корпоративный маркетинговый сайт сегодня является не
только представительской витриной, но и рабочим инструментом бизнеса. Через него
публикуются контентные страницы, статьи, проекты, вакансии, формы захвата лидов и другие
элементы, влияющие на маркетинговую активность компании.

На практике такие сайты часто развиваются быстрее, чем их контентная архитектура. В
результате часть данных остается зашитой во frontend-коде, публикация изменений требует
участия разработчика, а локализация, предпросмотр и \texttt{SEO}-настройки оказываются
фрагментированными.

\textbf{Объект исследования:} процессы управления контентом корпоративного маркетингового
сайта.

\textbf{Предмет исследования:} архитектурные и программные решения для построения
\texttt{CMS-first} веб-платформы на базе \texttt{Strapi} и \texttt{Astro}.

\section*{2. Краткая характеристика работы}

\begin{longtable}{|p{5.0cm}|p{9.0cm}|}
\hline
\textbf{Показатель} & \textbf{Значение} \\
\hline
Тип разрабатываемого решения & \texttt{CMS-first} платформа для корпоративного маркетингового сайта \\
\hline
Основной backend стек & \texttt{Strapi 5}, \texttt{SQLite}, \texttt{OpenAPI}, \texttt{Docker} \\
\hline
Основной frontend стек & \texttt{Astro 6}, \texttt{React}, \texttt{Tailwind}, \texttt{Node adapter}, \texttt{Docker}, \texttt{sitemap} \\
\hline
Архитектурная схема & \texttt{Strapi -> API -> Astro} \\
\hline
Основные контентные сущности & \texttt{global}, \texttt{home-page}, \texttt{page}, \texttt{article}, \texttt{author}, \texttt{project}, \texttt{vacancy}, \texttt{industry}, \texttt{job-role} \\
\hline
Прикладные сущности & \texttt{lead-submission}, \texttt{vacancy-application} \\
\hline
Основной функциональный результат & Управляемая контентная модель, публичная витрина, \texttt{preview}, \texttt{SEO}, формы и публикационный контур \\
\hline
Публикационный подход & \texttt{publish -> webhook -> rebuild} \\
\hline
\end{longtable}

\section*{3. Цель и задачи работы}

\textbf{Цель работы} --- разработать систему управления контентом для корпоративного
маркетингового сайта на базе \texttt{Strapi} и \texttt{Astro}, обеспечивающую управление
страницами, статьями, проектами и вакансиями, поддержку локалей \texttt{ru/en},
предварительного просмотра, \texttt{SEO}-контур и воспроизводимый сценарий публикации.

Для достижения цели решались следующие задачи:
\begin{enumerate}
  \item проанализировать особенности маркетинговых сайтов и обосновать выбор
  \texttt{CMS-first} архитектуры;
  \item спроектировать модель данных и маршрутный контур системы;
  \item реализовать управление страницами и ключевыми контентными сущностями через
  \texttt{Strapi};
  \item реализовать публичную витрину на \texttt{Astro} с локализацией, \texttt{preview},
  \texttt{SEO} и \texttt{sitemap};
  \item организовать публикационный и эксплуатационный контур, включая
  \texttt{webhook -> rebuild};
  \item провести проверку ключевых пользовательских и редакторских сценариев.
\end{enumerate}

\section*{4. Архитектурная схема и логика системы}

В работе реализован монорепозиторий на базе \texttt{Nx + pnpm}, включающий два основных
приложения:
\begin{itemize}
  \item \textbf{CMS-слой} на \texttt{Strapi 5}, отвечающий за хранение и редактирование
  контента;
  \item \textbf{frontend-слой} на \texttt{Astro 6}, отвечающий за публикацию контента на
  публичной витрине.
\end{itemize}

Контентная модель включает:
\begin{itemize}
  \item глобальные данные сайта: \texttt{global};
  \item главную страницу: \texttt{home-page};
  \item локализуемые CMS-страницы: \texttt{page};
  \item контентные коллекции: \texttt{article}, \texttt{author}, \texttt{project};
  \item карьерный модуль: \texttt{vacancy}, \texttt{industry}, \texttt{job-role};
  \item прикладные формы: \texttt{lead-submission}, \texttt{vacancy-application}.
\end{itemize}

Для \texttt{home-page} и \texttt{page} реализован конструктор страниц на основе
\texttt{Dynamic Zone}, который позволяет собирать страницу из переиспользуемых блоков без
правки frontend-кода.

На стороне frontend реализован locale-prefixed маршрутный контур:
\begin{itemize}
  \item \texttt{/:locale/};
  \item \texttt{/:locale/:slug/};
  \item \texttt{/:locale/articles/...};
  \item \texttt{/:locale/projects/...};
  \item отдельный карьерный контур \texttt{/vacancies/...};
  \item серверные preview-маршруты \texttt{/preview/...}.
\end{itemize}

\section*{5. Реализованный функциональный результат}

По итогам реализации подтверждены следующие результаты:
\begin{enumerate}
  \item реализовано управление главной страницей, CMS-страницами, статьями, проектами,
  вакансиями, маркетинговыми лидами и откликами через \texttt{Strapi};
  \item реализован безопасный \texttt{preview mode} для \texttt{home-page}, \texttt{page},
  \texttt{article}, \texttt{project}, \texttt{vacancy};
  \item реализован \texttt{SEO/Open Graph}-контур для ключевых сущностей и генерация
  \texttt{sitemap};
  \item реализованы server-side маршруты для форм \texttt{lead-submission} и
  \texttt{vacancy-application};
  \item реализован публикационный контур \texttt{publish -> webhook -> rebuild};
  \item реализована ролевая модель с разделением маркетингового и карьерного контуров.
\end{enumerate}

\section*{6. Контуры доступа и роли}

В системе выделены следующие роли:
\begin{itemize}
  \item \texttt{administrator} --- полный административный доступ;
  \item \texttt{marketer/content-manager} --- публикация и управление маркетинговым контентом;
  \item \texttt{editor} --- подготовка маркетинговых черновиков без права публикации;
  \item \texttt{hr} --- управление вакансиями и обработка откликов.
\end{itemize}

Это позволяет разделить редакторские процессы:
\begin{itemize}
  \item маркетинговый контур: \texttt{global}, \texttt{home-page}, \texttt{page},
  \texttt{article}, \texttt{author}, \texttt{project}, \texttt{lead-submission};
  \item карьерный контур: \texttt{vacancy}, \texttt{industry}, \texttt{job-role},
  \texttt{vacancy-application}.
\end{itemize}

\section*{7. Проверка сценариев и подтвержденные результаты}

Для проекта подготовлена воспроизводимая матрица приемки. В контрольной серии проверок
подтверждено:
\begin{itemize}
  \item публичные маршруты storefront-core открываются корректно;
  \item \texttt{preview} с неверным \texttt{secret} возвращает \texttt{401}, а
  draft-preview отдает страницу с \texttt{noindex};
  \item формы проходят server-side валидацию;
  \item валидные отправки форм сохраняются в \texttt{SQLite};
  \item frontend-сборка проходит успешно и формирует \texttt{37} prerendered static routes;
  \item генерируются \texttt{sitemap-index.xml} и \texttt{sitemap-0.xml};
  \item в базе данных подтвержден активный \texttt{Frontend rebuild hook} на
  \texttt{entry.publish} и \texttt{entry.unpublish};
  \item размер клиентского build output составляет около \texttt{2.4 МБ}.
\end{itemize}

\section*{8. Практическая значимость}

Практическая значимость работы заключается в переносе типовых редакторских операций из
frontend-кода в \texttt{CMS}. Контент-менеджер получает возможность управлять страницами,
статьями, проектами и метаданными через административный интерфейс, а разработчик
перестает быть обязательным участником каждого небольшого изменения контента.

Для проекта в целом это дает:
\begin{itemize}
  \item централизованную модель данных;
  \item единый публикационный контур;
  \item поддержку локализации;
  \item управляемый \texttt{preview};
  \item основу для дальнейшего масштабируемого развития публичной витрины.
\end{itemize}

\section*{9. Направления развития}

В ходе выполнения работы сформированы следующие направления дальнейшего развития:
\begin{enumerate}
  \item расширение англоязычного наполнения detail-разделов
  \texttt{articles/projects};
  \item дальнейшее развитие \texttt{CMS SEO}-контура и углубление редакторской настройки
  публичных страниц;
  \item расширение browser-level контура оценки доступности и производительности.
\end{enumerate}

\section*{10. Основные выводы}

\begin{enumerate}
  \item в результате выполнения выпускной квалификационной работы разработана рабочая
  \texttt{CMS-first} система управления контентом для корпоративного маркетингового сайта;
  \item полученное решение объединяет контентную модель, публичную витрину, локализацию,
  \texttt{preview}, формы взаимодействия с пользователем и публикационный контур;
  \item практический результат работы состоит в переходе от кодо-центричного
  подхода к централизованной модели управления контентом;
  \item разработанная система может рассматриваться как инженерная основа для дальнейшего
  развития маркетинговой платформы как управляемой и воспроизводимой системы.
\end{enumerate}

\end{document}
